\scanpage{125}
On déduit de (194) un morphisme correspondant (qui est épi lisse)

(197)\quad
\diagram{\begin{tikzcd}[ampersand replacement=\&,column sep=huge,row sep=large]
G(d,d')_d\arrow[r]\arrow[d,equal,"\mathrm{def}"]\&
G(d,d')_{\unsure{d'|d''}}\arrow[d,equal,"\mathrm{def}"]\\
\begin{gathered}G_d\mathbin\times_{D_d}D_{d,d'}\\
\simeq T(d,d')_d\otimes_{\mathbb Z}\mathbb G_m\end{gathered}\&
\begin{gathered}G_{d'}\mathbin\times_{D_{d'}}D_{d,d'}\\
\simeq T(d,d')_{d'}\otimes_{\mathbb Z}\mathbb G_m\end{gathered}
\end{tikzcd}}
et on \unsure{voit|sait} que l'on a un isom. correspondant

(198)\quad
\diagram{\begin{tikzcd}[ampersand replacement=\&,column sep=small,row sep=large]
P(d,d')_{\unsure{d'|d''}}\arrow[d,equal,"\mathrm{def}"]\&
\simeq\&
\underbrace{P(d,d')_d}\arrow[d,equal,"\mathrm{def}"]\&
{}\mathbin\wedge_{\unsure{G_d|G}(d,d')_d}G(d,d')_{d'}\\
P_d\mathbin{\unsure{\otimes|\times}}_{D_d}D_{d,d'}\&\&
P_{d'}\mathbin\times_{D_{d'}}D_{d,d'}\&
\end{tikzcd}}

\inserted{On a des relations de transitivité}
\deleted{Au niveau des groupoïdes fondamentaux, on peut dire que $\Pi P(d,d')_{d'}$}
\inserted{pour les morphismes précédents,}
\deleted{$\unsure{\Pi P(d,d')_d}$}
\unsure{relatifs|relatives} a donnée d'un triple
\[
d\geq d'\geq d''
\]
en se plaçant sur $D(d,d',d'')$, où on a des \deleted{objets} systèmes locaux
\[
(199)\qquad T(d,d',d'')_d,\quad T(d,d',d'')_{d'},\quad T(d,d',d'')_{d''},
\]
des tores correspondants
\[
(200)\qquad G(d,d',d'')_d,\quad G(d.d',d'')_{d'},\quad G(d,d',d'')_{d''}
\]
\margin{124}
